\documentclass[11pt, a4paper]{article}

\input{bfw-style.tex}

\title{\vspace*{-1.5cm}\Huge\bfseries\color{bfwPrimaryDark}Aufgabenheft Session~4\\[0.4em]
  \large\color{bfwInkSoft}\textnormal{Rechtsformen \& Aufbauorganisation --- Übungen im IHK-Klausurformat}}
\author{\textsf{IT Lernfeld 1 --- Das Unternehmen und die eigene Rolle im Betrieb beschreiben}}
\date{Selbstlernmaterial}

\begin{document}
\maketitle
\thispagestyle{empty}

\begin{merksatz}[title={\bfseries So nutzen Sie dieses Aufgabenheft}]
Bearbeiten Sie alle Aufgaben zunächst \emph{ohne} in die Lösungen zu schauen. Schreiben
Sie Ihre Antworten in vollständigen Sätzen --- die IHK-Prüfung verlangt formulierte
Begründungen, nicht bloße Stichworte. Beachten Sie die \emph{Operatoren} (nennen,
erläutern, beschreiben, begründen, beurteilen) --- sie geben den geforderten
Antwort-Tiefegrad vor. Erst danach öffnen Sie den Lösungsteil ab Seite~\pageref{loesungen}
und vergleichen.
\end{merksatz}

\bigskip

\textbf{Kontext:} Sie arbeiten seit Kurzem als Auszubildende bei der \textbf{JIKU IT-Solutions GmbH}
in Hamburg. Im Rahmen Ihrer Ausbildung beschäftigen Sie sich heute mit der rechtlichen Struktur
Ihres Ausbildungsbetriebs und der internen Aufbauorganisation. Ihre Ausbilderin Frau Meyerhoff
begleitet Sie dabei.

\bigskip

\noindent
\begin{tabular}{|l|l|l|l|l|l|l|l|}
\hline
\textbf{Aufgabe} & 1 & 2 & 3 & 4 & 5 & 6 & \textbf{Gesamt}\\
\hline
\textbf{Punkte} & /10 & /8 & /8 & /10 & /10 & /10 & /56\\
\hline
\end{tabular}

\tableofcontents
\vspace{1em}
\hrule

\section{Aufgaben}

\subsection*{Aufgabe~1 --- Rechtsformvergleich (10 Punkte)}

\begin{enumerate}[label=(\alph*)]
  \item \textbf{Nennen} Sie die drei Kapitalgesellschaften, die in der deutschen Rechtspraxis besonders
  verbreitet sind, und geben Sie jeweils das gesetzliche Mindestkapital an. \hfill (6~P.)

  \vspace{5cm}

  \item \textbf{Erläutern} Sie den Unterschied zwischen der Haftung bei einer OHG und der Haftung
  bei einer GmbH. \hfill (4~P.)

  \vspace{4.5cm}
\end{enumerate}

\subsection*{Aufgabe~2 --- Kaufmannseigenschaft und Handelsregister (8 Punkte)}

\begin{enumerate}[label=(\alph*)]
  \item \textbf{Beschreiben} Sie, welche Funktion das Handelsregister in der deutschen Wirtschaft hat,
  und \textbf{nennen} Sie die beiden Abteilungen mit je einem Beispiel. \hfill (4~P.)

  \vspace{5cm}

  \item JIKU IT-Solutions GmbH ist im Handelsregister eingetragen. \textbf{Erläutern} Sie, welche
  rechtlichen Konsequenzen die Kaufmannseigenschaft für JIKU hat. Nennen Sie zwei konkrete
  Pflichten. \hfill (4~P.)

  \vspace{4cm}
\end{enumerate}

\subsection*{Aufgabe~3 --- Firma und Firmengrundsätze (8 Punkte)}

\begin{enumerate}[label=(\alph*)]
  \item \textbf{Erklären} Sie, was man im Handelsrecht unter der „Firma" eines Unternehmens versteht,
  und \textbf{grenzen} Sie diesen Begriff vom alltäglichen Sprachgebrauch ab. \hfill (2~P.)

  \vspace{3cm}

  \item \textbf{Ordnen} Sie die folgenden Firmennamen den drei Firmenarten zu (Personenfirma, Sachfirma
  oder Fantasiefirma) und \textbf{begründen} Sie kurz. \hfill (3~P.)

  \bigskip

  \noindent
  \begin{tabular}{|p{4.5cm}|p{3.5cm}|p{5cm}|}
  \hline
  \textbf{Firmenname} & \textbf{Firmenart} & \textbf{Begründung}\\
  \hline
  IT-Sander e.\,K. & & \\[0.8cm]
  \hline
  Lasse Müller e.\,Kfm. & & \\[0.8cm]
  \hline
  Sysygy e.\,Kfm. & & \\[0.8cm]
  \hline
  \end{tabular}

  \bigskip

  \item \textbf{Nennen} und \textbf{erläutern} Sie die drei Firmengrundsätze nach dem HGB. \hfill (3~P.)

  \vspace{5cm}
\end{enumerate}

\subsection*{Aufgabe~4 --- Prokura und Handlungsvollmacht (10 Punkte)}

Frau Meyerhoff erklärt Ihnen, dass Mitarbeiterinnen und Mitarbeiter je nach Aufgabe unterschiedliche
Vollmachten erhalten.

\begin{enumerate}[label=(\alph*)]
  \item \textbf{Unterscheiden} Sie Prokura und Handlungsvollmacht hinsichtlich Umfang, Zeichnung
  und Erteilung. \hfill (6~P.)

  \vspace{7cm}

  \item Der Buchhalter Herr Garza unterschreibt einen Liefervertrag mit dem Zusatz „i.\,A.\,Garza".
  Die Einkäuferin Frau Pelz zeichnet einen Rahmenvertrag mit „i.\,V.\,Pelz". Der Prokurist Herr Volk
  unterschreibt mit „ppa.\,Volk". \textbf{Ordnen} Sie jedem die korrekte Vollmachtsart zu und
  \textbf{begründen} Sie kurz. \hfill (4~P.)

  \vspace{5cm}
\end{enumerate}

\subsection*{Aufgabe~5 --- Aufbauorganisation und Leitungssysteme (10 Punkte)}

\begin{enumerate}[label=(\alph*)]
  \item \textbf{Unterscheiden} Sie Aufbauorganisation und Ablauforganisation. Geben Sie je ein
  konkretes Beispiel aus JIKU IT-Solutions. \hfill (4~P.)

  \vspace{4.5cm}

  \item JIKU hat laut Organigramm Stabsstellen eingerichtet. \textbf{Erläutern} Sie, was eine
  Stabsstelle ist, und nennen Sie zwei Merkmale, die sie von einer Instanz unterscheiden. \hfill (3~P.)

  \vspace{4cm}

  \item \textbf{Beschreiben} Sie das Leitungssystem, das JIKU IT-Solutions nach dem Lehrbuch
  eingerichtet hat, und \textbf{begründen} Sie, warum es für das IT-Systemhaus geeignet ist. \hfill (3~P.)

  \vspace{4cm}
\end{enumerate}

\subsection*{Aufgabe~6 --- Rechtsformwahl (10 Punkte)}

Frau Meyerhoff stellt Ihnen drei fiktive Gründungsvorhaben vor:

\begin{itemize}
  \item \textbf{Vorhaben~A:} Zwei IT-Fachinformatikerinnen wollen gemeinsam Webentwicklungs-Projekte
  übernehmen. Sie haben kein Startkapital, wollen aber schnell loslegen und keine aufwendige
  Gründungsformalität.
  \item \textbf{Vorhaben~B:} Ein IT-Gründer möchte ein innovatives Start-up für KI-gestützte
  Softwarelösungen aufbauen, haftet aber grundsätzlich nicht mit seinem Privatvermögen. Er hat
  8.000~€ Eigenkapital.
  \item \textbf{Vorhaben~C:} Ein etabliertes IT-Systemhaus will weiterwachsen, externe Investoren
  anziehen und Mitarbeitende durch Unternehmensanteile am Erfolg beteiligen.
\end{itemize}

\begin{enumerate}[label=(\alph*)]
  \item \textbf{Empfehlen} Sie für jedes Vorhaben eine passende Rechtsform und \textbf{begründen}
  Sie Ihre Entscheidung anhand von mindestens zwei Kriterien. \hfill (6~P.)

  \vspace{7cm}

  \item \textbf{Erläutern} Sie den steuerlichen Unterschied zwischen Personen- und Kapitalgesellschaften
  und nennen Sie die jeweilige Steuerart. \hfill (4~P.)

  \vspace{4.5cm}
\end{enumerate}

% =====================================================================
\newpage
\section*{Lösungshinweise für Lehrende}
\label{loesungen}
\addcontentsline{toc}{section}{Lösungshinweise für Lehrende}
% =====================================================================

\begin{merksatz}[title={\bfseries Hinweis zur Nutzung}]
Dieser Abschnitt ist ausschließlich für Lehrende bestimmt. Die Lösungen geben
Musterlösungen und typische Fehlerquellen an. Teilnehmende erhalten diesen Teil
erst nach der eigenständigen Bearbeitung.
\end{merksatz}

\subsection*{Lösung Aufgabe~1}

\begin{enumerate}[label=(\alph*)]
  \item Drei Kapitalgesellschaften:
  \begin{itemize}
    \item \textbf{GmbH (Gesellschaft mit beschränkter Haftung):} Mindest-Stammkapital
    \textbf{25.000~€} (§~5 GmbHG); bei Gründung mindestens 12.500~€ sofort einzuzahlen.
    \item \textbf{UG (haftungsbeschränkt):} Mindest-Stammkapital \textbf{ab 1~€}; Thesaurierungspflicht:
    25\,\% des Jahresüberschusses zurücklegen, bis 25.000~€ erreicht sind.
    \item \textbf{AG (Aktiengesellschaft):} Mindest-Grundkapital \textbf{50.000~€} (§~7 AktG).
  \end{itemize}
  \textit{Häufiger Fehler: GmbH-Mindestkapital (25.000~€) mit AG-Mindestkapital (50.000~€) verwechseln.}

  \item Bei der \textbf{OHG} haften alle Gesellschafter \emph{persönlich, unbeschränkt und solidarisch}
  --- jeder Gesellschafter haftet für die gesamten Schulden der Gesellschaft, auch mit seinem
  Privatvermögen. Bei der \textbf{GmbH} haftet \emph{ausschließlich die Gesellschaft} mit ihrem
  Gesellschaftsvermögen; Gesellschafter haften nicht persönlich (Haftungsbeschränkung).
\end{enumerate}

\subsection*{Lösung Aufgabe~2}

\begin{enumerate}[label=(\alph*)]
  \item Das Handelsregister beim Amtsgericht ist eine zentrale Informationsquelle für die Wirtschaft
  (Publikationsfunktion: jedermann kann Einsicht nehmen; Schutzfunktion: schützt Kunden und Lieferanten).
  \textbf{Abteilung~A:} Einzelkaufleute und Personengesellschaften (z.\,B.\ JIKU Hamburg e.\,K.,
  Netzwerk~OHG). \textbf{Abteilung~B:} Kapitalgesellschaften (z.\,B.\ JIKU IT-Solutions GmbH, DataCenter~AG).

  \item Als Kaufmann unterliegt JIKU neben dem BGB auch dem HGB. Zwei Pflichten (Auswahl):
  \textbf{doppelte Buchführung} einrichten (§~238 HGB), \textbf{Jahresabschluss} erstellen und
  offenlegen, Firma mit Rechtsformzusatz führen, Prokura erteilen können.
  \textit{Hinweis: Die Kaufmannseigenschaft entsteht durch Eintragung oder durch Betreiben eines
  kaufmännisch eingerichteten Handelsgewerbes.}
\end{enumerate}

\subsection*{Lösung Aufgabe~3}

\begin{enumerate}[label=(\alph*)]
  \item Die \textbf{Firma} ist nach § 17 HGB der Name, unter dem ein Kaufmann seine Geschäfte
  betreibt und die Unterschrift abgibt. Im alltäglichen Sprachgebrauch meint „Firma" oft das
  Unternehmen selbst --- juristisch bezeichnet sie nur den \emph{Handelsnamen} des Kaufmanns.
  Nur eingetragene Kaufleute dürfen eine Firma führen.

  \item Firmenarten:
  \begin{itemize}
    \item IT-Sander e.\,K.~: \textbf{Sachfirma} (Gegenstand der Tätigkeit).
    \item Lasse Müller e.\,Kfm.~: \textbf{Personenfirma} (Familienname des Inhabers).
    \item Sysygy e.\,Kfm.~: \textbf{Fantasiefirma} (freie Namensgestaltung ohne sachlichen Bezug).
  \end{itemize}

  \item Firmengrundsätze:
  \begin{itemize}
    \item \textbf{Firmenausschließlichkeit} (§~30 HGB): Die Firma muss sich von anderen Firmen
    am selben Ort klar unterscheiden --- keine Verwechslungsgefahr.
    \item \textbf{Firmenwahrheit} (§~18 Abs.~2 HGB): Irreführungsverbot --- keine täuschenden
    Angaben über Art, Umfang oder sonstige Verhältnisse des Unternehmens.
    \item \textbf{Firmenbeständigkeit} (§~22 HGB): Die Firma kann auch bei Inhaberwechsel
    fortgeführt werden; sie ist durch § 12 BGB, § 37 HGB und § 15 MarkenG geschützt.
  \end{itemize}
\end{enumerate}

\subsection*{Lösung Aufgabe~4}

\begin{enumerate}[label=(\alph*)]
  \item

  \begin{tabular}{|p{3cm}|p{5.5cm}|p{5cm}|}
  \hline
  \textbf{Merkmal} & \textbf{Prokura (§§ 48--53 HGB)} & \textbf{Handlungsvollmacht (§§ 54--58 HGB)}\\
  \hline
  Umfang & Alle gewöhnlichen und außergewöhnlichen Geschäfte; ausgenommen: Bilanz/Steuererklärung, HR-Eintragungen, Prokura erteilen & Gewöhnliche Geschäfte (Gesamt), bestimmte Art (Art), Einzelgeschäft (Einzel)\\
  \hline
  Zeichnung & ppa.~Name (per procura) & i.V.~Name oder i.A.~Name \newline {\small\itshape (kaufmännische Konvention, gesetzlich nicht vorgeschrieben)}\\
  \hline
  Erteilung & Nur durch Kaufmann; Eintragung ins HR vorgeschrieben (deklaratorisch) & Formlos (mündlich/schriftlich); kein HR-Eintrag möglich\\
  \hline
  \end{tabular}

  \item
  \begin{itemize}
    \item Herr Garza (i.\,A.): \textbf{Einzelvollmacht} --- „i.\,A." steht für „im Auftrag"
    und steht für die Einzelvollmacht (i.\,V.\ wird bei Gesamt- und Artvollmacht verwendet ---
    kaufmännische Konvention).
    \item Frau Pelz (i.\,V.): \textbf{Gesamtvollmacht (allgemeine Handlungsvollmacht)} ---
    „i.\,V." steht für „in Vollmacht", typisch für die allgemeine Handlungsvollmacht.
    \item Herr Volk (ppa.): \textbf{Prokura} --- „ppa." (per procura) ist der gesetzlich
    vorgeschriebene Zeichnungszusatz für Prokuristen.
  \end{itemize}
\end{enumerate}

\subsection*{Lösung Aufgabe~5}

\begin{enumerate}[label=(\alph*)]
  \item \textbf{Aufbauorganisation:} Gliederung des Betriebes in Einheiten, Übertragung von Aufgaben
  und Kompetenzen (z.\,B.\ Organigramm mit Abteilungen IT-Service, Cloud-Hosting, Verwaltung bei
  JIKU). \textbf{Ablauforganisation:} Regelung des Arbeitsablaufs in zeitlicher, räumlicher und
  funktionaler Hinsicht (z.\,B.\ Ticketbearbeitungsprozess im Help-Desk: Anfrage --- Diagnose ---
  Lösung --- Rückmeldung).

  \item Eine \textbf{Stabsstelle} ist eine Hilfsstelle einer Instanz, die berät und Entscheidungen
  vorbereitet, aber \emph{keine Weisungsbefugnis} gegenüber anderen Stellen hat. Zwei Unterschiede
  zur Instanz: (1) Instanz hat Anordnungsbefugnis, Stabsstelle nicht. (2) Stabsstelle hat kein
  Weisungsrecht, aber ein Informationsrecht aus allen Linienstellen.

  \item JIKU hat laut Lehrbuch eine \textbf{Kombination aus Stabliniensystem und Matrixorganisation}
  eingerichtet. Das Stabliniensystem sorgt für klare Linienkommunikation und Entlastung durch
  Stabsstellen. Die Matrixorganisation ist für die kundennahen IT-Solutions-Bereiche geeignet, in
  denen Projektarbeit mit unterschiedlichen Fachkompetenzen überwiegt.
\end{enumerate}

\subsection*{Lösung Aufgabe~6}

\begin{enumerate}[label=(\alph*)]
  \item Empfehlungen (Auswahl; gut begründete Alternativen anerkennbar):
  \begin{itemize}
    \item \textbf{Vorhaben~A:} \textbf{GbR} --- einfachste Mehrpersonengesellschaft, formlos per
    Vertrag gründbar, kein Mindestkapital. Nachteil: persönliche Haftung, keine eigene Firma.
    Alternative gut begründbar: OHG (wenn Handelsgewerbe beabsichtigt).
    \item \textbf{Vorhaben~B:} \textbf{UG (haftungsbeschränkt)} --- ab 1~€ Stammkapital; Haftung
    beschränkt auf Gesellschaftsvermögen; bei 25.000~€ Aufstieg zur GmbH möglich. Kriterien:
    Haftungsbeschränkung, geringer Kapitalbedarf.
    \item \textbf{Vorhaben~C:} \textbf{AG} --- einfache Beteiligung externer Investoren und
    Mitarbeitender über Aktien; Kapital beschaffbar über Börse. Kriterien: Kapitalaufnahme,
    Mitarbeiterbeteiligung. Mindestgrundkapital: 50.000~€.
  \end{itemize}
  \textit{Häufiger Fehler: UG und GmbH verwechseln; die UG hat kein Mindestkapital von 25.000~€,
  sondern ab 1~€ mit Thesaurierungspflicht.}

  \item \textbf{Personengesellschaften} (Einzelunternehmen, GbR, OHG, KG): Jeder Gesellschafter
  ist selbst Steuersubjekt und zahlt \textbf{Einkommensteuer} (§§~15, 18 EStG) auf seinen
  Gewinnanteil. \textbf{Kapitalgesellschaften} (GmbH, AG): Die Gesellschaft ist als juristische
  Person eigenes Steuersubjekt und zahlt \textbf{Körperschaftsteuer} (§~1 KStG) auf ihren Gewinn.
\end{enumerate}

\end{document}
